% Target: rmp-ii-mpu-normal-form
% Formula: A virtual MPO line pulls through a PEPS tensor without changing the state.
% Ink: Kernel plane frames carry the PEPS stars and marked virtual paths.
% Boundary: Each PEPS atom has four virtual and one physical open leg; each marked
%   virtual path has two open endpoints.
% Source: paper TN-Review-main.tex line 404; author source ImagesReview Section
%   II/ImagesReview Section II.tex lines 339-357
% Capabilities: kernel, plane-frame, typed-ports, open-string
\[
  \tenkzkernel
  \tndeclare{species}{gauge}{hue=source:blue}
  \begin{tenkz}[
      rows={wire,wire,wire}, cols=3, frame=plane, bonds=none, physical=up]
    \tn[at=(2,2), name=A,
      ports={0:virtual,90:virtual,180:virtual,270:virtual}]{}
    \tnwire{A.180}{open w} \tnwire{A.0}{open e}
    \tnwire{A.90}{open n} \tnwire{A.270}{open s}
    \tnwire[kind=string, species=gauge, stroke=dotted, name=gA,
      via={(2,1),(3,2)}]{0.5 n of (2,1)}{0.5 s of (3,2)}
    \tn[at=(2,1), physical=none, skin=dot, species=gauge]{}
    \tn[at=(3,2), physical=none, skin=dot, species=gauge]{}
  \end{tenkz}
  \;=\;
  \begin{tenkz}[
      rows={wire,wire,wire}, cols=3, frame=plane, bonds=none, physical=up]
    \tn[at=(2,2), name=B,
      ports={0:virtual,90:virtual,180:virtual,270:virtual}]{}
    \tnwire{B.180}{open w} \tnwire{B.0}{open e}
    \tnwire{B.90}{open n} \tnwire{B.270}{open s}
    \tnwire[kind=string, species=gauge, stroke=dotted, name=gB,
      via={(2,3),(1,2)}]{0.5 s of (2,3)}{0.5 n of (1,2)}
    \tn[at=(2,3), physical=none, skin=dot, species=gauge]{}
    \tn[at=(1,2), physical=none, skin=dot, species=gauge]{}
  \end{tenkz}
\]
